\chapter{Nipple Discharge}

\section*{Definition}
Any fluid that escapes from the nipple, varying in character (milky, serous, purulent, bloody, green/black) and significance from benign physiological discharge to a marker of underlying breast pathology including infection, duct ectasia, papilloma, or malignancy.

\section*{What Happens in Your Body}
Milk production is regulated by prolactin from the anterior pituitary, normally suppressed by dopamine from the hypothalamus. Pathologic discharge arises from spontaneous activation of breast epithelial cells (hyperprolactinemia), ductal inflammation (duct ectasia, periductal mastitis), intraductal growth (papilloma), or malignant cells eroding into ducts. Unilateral, spontaneous, bloody discharge from a single duct raises highest concern for malignancy.

\section*{Causes}
\begin{itemize}
  \item Physiological discharge (bilateral, milky, expressed only, common in reproductive-age women)
  \item Intraductal papilloma (most common cause of bloody unilateral discharge)
  \item Duct ectasia (benign dilation of mammary ducts, thick green/black discharge)
  \item Mastitis or breast abscess (purulent discharge with pain and erythema)
  \item Hyperprolactinemia (pituitary adenoma, medications, hypothyroidism)
  \item Breast cancer (invasive or DCIS, often unilateral spontaneous bloody discharge)
  \item Pregnancy and lactation (colostrum, milk)
  \item Galactorrhea (medication-induced: antipsychotics, metoclopramide, SSRIs, oral contraceptives)
  \item Breast trauma or surgery
  \item Fibrocystic breast changes
\end{itemize}

\section*{When to See a Doctor}
See your doctor if:
\begin{itemize}
  \item Symptoms are severe or getting worse over time
  \item Spontaneous, unilateral, bloody, or clear watery discharge; discharge associated with a palpable breast lump; or persistent discharge requiring pad use
  \item You have other concerning symptoms such as fever, weight loss, or bleeding
\end{itemize}

\section*{Related Symptoms}
\begin{itemize}
  \item Breast lump
  \item Breast pain
  \item Breast swelling
  \item Galactorrhea
  \item Nipple changes
\end{itemize}
\textbf{Important:} If this symptom persists, your doctor will check for serious underlying causes.

\section*{Self-Care / Home Management}
\begin{itemize}
  \item Do not squeeze or attempt to express nipple discharge repeatedly, as this can stimulate further secretion and cause local irritation.
  \item Keep a record of the discharge characteristics: colour (milky, clear, yellow, green, bloody), whether it is spontaneous or expressed, unilateral or bilateral, and which duct it comes from.
  \item Wear a cotton bra with disposable breast pads if the discharge is bothersome or stains clothing.
  \item Avoid nipple stimulation and check for any associated breast lump, skin changes, or breast tenderness that should be reported to your doctor.
\end{itemize}

\section*{How Your Doctor Will Evaluate You}
\begin{itemize}
  \item Clinical breast examination assessing whether the discharge is spontaneous or expressed, unilateral or bilateral, single-duct or multi-duct, and the colour/consistency.
  \item Mammography (with tomosynthesis) and targeted ultrasound are first-line imaging to evaluate for intraductal papilloma, duct ectasia, or malignancy.
  \item Cytological analysis of discharge fluid may be performed, though its sensitivity for malignancy is limited; ductography (galactogram) can identify intraductal lesions.
  \item MRI breast with contrast is indicated when there is suspicious discharge (spontaneous, unilateral, bloody) with negative mammogram and ultrasound.
\end{itemize}

\section*{Treatment Options}
\begin{itemize}
  \item Physiological discharge requires no treatment beyond reassurance and avoidance of nipple stimulation.
  \item Intraductal papilloma is managed with surgical excision (microdochectomy) to relieve symptoms and confirm benign histology.
  \item Duct ectasia is treated conservatively with warm compresses and NSAIDs; surgical excision of the affected duct is reserved for persistent or bothersome discharge.
  \item Malignant discharge (associated with DCIS or invasive cancer) requires surgical excision (wide local excision or mastectomy) and adjuvant therapy as per breast cancer guidelines.
  \item Medications, lifestyle modifications, and physical therapies may be prescribed as appropriate.
  \item Follow-up ensures treatment effectiveness and allows timely adjustments to the care plan.
\end{itemize}

\section*{References}
\begin{thebibliography}{99}
\bibitem{nippledisc1} Expert Panel on Breast Imaging, Lee SJ, Trikha S, et al. "ACR Appropriateness Criteria evaluation of nipple discharge." J Am Coll Radiol. 2020;17(5S):S138-S147.
\bibitem{nippledisc2} Seltzer MH, Perloff LJ, Kelley RI, et al. "The significance of nipple discharge." Am Surg. 1970;36(1):17-22.
\bibitem{nippledisc3} Tabar L, Dean PB, Pentek Z. "Galactography in the diagnosis of nipple discharge." Radiology. 2017;282(2):351-359.
\bibitem{nippledisc4} Boccardo F, Puntoni M, Menichini E, et al. "Nipple discharge: a review of diagnostic modalities." Breast J. 2021;27(1):54-61.
\bibitem{nippledisc5} Smith RA, Andrews KS, Brooks D, et al. "Cancer screening in the United States, 2019." CA Cancer J Clin. 2019;69(3):184-210.
\bibitem{nippledisc6} Harris JR, Lippman ME, Morrow M, et al. "Diseases of the breast." 6th ed. Wolters Kluwer; 2020.
\end{thebibliography}

\clearpage